\section{Cosmic Environment}

In \emph{The Great Triad,} \textbf{Rene Guenon} discusses the circumstances and environment in which a given being manifests. First of all, the cosmic environment is a unity, as he writes:

\begin{quotex}
The cosmic environment can only be conceived of as a whole of which each part is linked to every other part without any break in continuity. To try to conceive of it otherwise would be to assume the existence of a void, but this is not a possibility of manifestation and can have no place in the Cosmos.

\end{quotex} 

The following assertion is difficult to accept on strictly material and physical terms. The Self chooses the environment. This is radically contrary to the existentialists like Heidegger who claims that we are ``thrown into the world" somehow by accident and against our will. Guenon says the opposite:

\begin{quotex}
It is the individual being itself which by its own nature determines the specific conditions of its manifestation.

\end{quotex}
The entire modern world denies that there is any such nature; they call such a view ``essentialism". So birth, to them, is a matter of ``luck", something purely arbitrary. That is because they don't take into account that heredity is not just corporeal, but also subtle, as Guenon explains:

\begin{quotex}
If we take the case of heredity, we can say that there is not only a physiological heredity, but also a psychic heredity. The two kinds of heredity have exactly the same explanation, both being due to the presence in a person's make-up of elements derived from the particular environment into which the individual was born. In the case of the individual human state, it is obvious that the environmental elements will belong to the different modalities of that specific state: that is, some will be corporeal, others subtle or psychic. A person's make-up of elements derived from the particular environment into which the individual was born. 

\end{quotex}
These modalities, then, will include the following elements:

\begin{itemize}
\item \textbf{Physical}: This includes the material elements of the environment, including DNA. 
\item \textbf{Etheric}: This level includes all the life forces in the environment. Together with the next level, they make up the biosphere. 
\item \textbf{Astral}: This includes all the elements of sentient life 
\item \textbf{Noetic}: This modality includes the intellectual currents of the environment. This is called the Noosphere. 
\end{itemize}
Nevertheless, these physiological and psychical influences do not determine the human being. Guenon makes this clear:

\begin{quotex}
The two terms represent only the superficial and external aspects of the manifested being and are no more than simple modalities of one and the same state of existence — the state of existence represented by the horizontal plane. In short, the physiological and the psychic are just as contingent as each other, and the true being is beyond them both. 

\end{quotex}
That has consequences. Every being is related to every other being, either diachronically or synchronically.

\begin{quotex}
In reality these influences extend much further; indeed we can say quite literally and without the slightest exaggeration that they extend indefinitely in every possible direction. In fact the \emph{cosmic environment}, which is the domain of the state of manifestation we are considering, can only be conceived of as a whole of which each part is linked to every other part without any break in continuity. To try to conceive of it otherwise would be to assume the existence of a `void'. but this is not a possibility of manifestation and can have no place in the Cosmos. It follows that there must necessarily be relationships — which is basically to say that there must be reciprocal actions and reactions — between all the individual beings which are manifested in this domain, either simultaneously or successively. 

\end{quotex}
\paragraph{Affinity}
This implies that the being, when entering the human state, chooses the environment that best suits its nature.

\begin{quotex}
A being's situation in the environment is in the last analysis determined by the nature of that particular being, and the elements which it takes from its immediate surroundings, as well as the elements — both subtle and corporeal — which it draws to itself from the indefinite totality of its particular domain of manifestation, will necessarily correspond in some way to that nature. 

\end{quotex}
And he continues:

\begin{quotex}
Here we have the explanation for that `affinity' which dictates that a being coming into manifestation will only take from the environment whatever conforms with the possibilities which it bears within itself, and which belong to it alone and to no other being. In other words, as a result of this conformity it will take whatever is necessary to provide the contingent conditions which will allow these possibilities to develop or `actualise' themselves in the course of its individual manifestation. 

\end{quotex}
\paragraph{Two Connected Beings}
The highest level of affinity is when two beings find that they belong to each other in a very special way. They can only actualize all their possibilities when they have support from each other. They are meant to be together from all eternity, no matter what obstacles are put in their paths.

\begin{quotex}
For any relationship between two beings to be real it must necessarily be the expression of something inherent in the natures of both. So, the influence that a being might appear to undergo from the outside and to receive from another being is never anything else, when looked at from a more profound point of view, than a sort of translation in relation to the environment of a possibility inherent in the nature of the being itself.

\end{quotex}
\paragraph{Summary}
To sum up: the being manifests at the point of the intersection of the vertical and horizontal axis, the former being the Self and the latter the environment in which he works out his possibilities of manifestation.

\begin{quotex}
In the first instance this individual nature proceeds or originates from what the being is in itself. In the second instance it proceeds from the sum total of the influences exerted on it by the environment in which it manifests. The first element represents the inward or active side of the being; the second its outward or passive side.

\end{quotex}


\flrightit{Posted on 2022-02-17 by Cologero }
